\section*{\texorpdfstring{\S\CurrentExerciseGroup}{第{\CurrentExerciseGroup}节}}
\phantomsection
\addcontentsline{toc}{section}{第{\CurrentExerciseGroup}节习题}

\begin{enumerate}
\def\labelenumi{\arabic{enumi})}
\item
  设第 1 号的定理 1 之假设成立。设 \(h\) 是局部 \(\mu\) -可积函数，使得 \(p\) 是 \((h \cdot \mu)\) -逆紧的。证明函数 \(b \mapsto g(b) = \int h(t) d\lambda_b(t)\) ——它对测度 \(\nu = p(\mu)\) 而言在 B 中局部几乎处处有定义——满足 \(p(h \cdot \mu) = g \cdot \nu\) 。
\item
  设 B 是局部紧空间，\((\nu_{\iota})_{\iota\in I}\) 是 B 上全体正测度构成的族。设 T 是乘积空间 \(I\times B\) （其中 I 赋有离散拓扑），设 \(\nu\) 是 T 上这样的测度：它在 \(\{\iota\}\times B\) 上的限制是 B 上测度 \(\nu_{\iota}\) 的典范像（对一切 \(\iota\in I\) ）。设 p 是 T 到 B 上的射影；证明：若 B 中存在不可数的紧子集，则 \(\nu\) 在 p 下不容许伪像测度。（证明对这样的测度，B 的每一点都将有测度 \textgreater0 。）
\item
  给出如下例子：正测度 \(\mu\) （在一个波兰局部紧空间 \(T\) 上的）与连续映射 \(p\) ( \(T\) 到一个波兰局部紧空间 \(B\) 中的)，使得 \(p\) 不是 \(\mu\) -逆紧的。
\item
  设 \(\mathrm{T}\) 是区间 {[}0,1{]} ( \(\mathbf{R}\) 中的)，\(\mu\) 是 \(\mathrm{T}\) 上的 Lebesgue 测度。设 \(\mathrm{R}\{x,y\}\) 为等价关系 \(x - y \in \mathbf{Q}\) （在 \(\mathrm{T}\) 中）。证明 \(\mathrm{R}\) 不是 \(\mu\) -可测的（应用第 4 号定理 3），但 \(\mathrm{R}\) 在 \(\mathrm{T} \times \mathrm{T}\) 中的图像对乘积测度 \(\mu \otimes \mu\) 是可忽略的。
\item
  设 T 是 Cantor 三分集 \(K \subset [0,1]\) 与区间 I = {[}1,2{]}（R 中的）的并，设 \(\mu\) 是 Lebesgue 测度在紧空间 T 上诱导的测度。设 P 是 I 的一个具有连续统势的不可测子集（第 IV 章 §4，习题 8），设 \(\psi\) 是 K 到 P 上的双射。考虑 T 中的等价关系 R：每个不属于 \(K \cup P\) 的点 x 自成一个等价类，而点 \(y \in K\) 的类由 y 与 \(\psi(y)\) 组成。证明 R 是 \(\mu\) -可测的，但 K 对 R 的饱和化不是 \(\mu\) -可测的。
\item
  设 T 是波兰局部紧空间，\(\mu\) 是 T 上的正测度，R 是 T 中的 \(\mu\) -可测等价关系，p 是 T 到波兰局部紧空间 B 中的 \(\mu\) -可测映射，使得 \(p(x) = p(y)\) 等价于 \(R\{x, y\}\) ，\(\nu\) 是 \(\mu\) 在 p 下的伪像测度，而 \(b \mapsto \lambda_b\) 是 \(\mu\) 按关系 R 的一种分解。对每个 \(b \in B\) (满足 \(\lambda_b \neq 0\) 的)，设 \(C(b)\) 是承载 \(\lambda_b\) 的模 R 等价类；证明 \(\varphi_{C(b)} \cdot \mu\) 与 \(\lambda_b\) 成比例。给出一个例子，其中 \(\varphi_{C(b)} \cdot \mu = 0\) 对每个 \(b \in B\) 成立。
\item
  设 T 是波兰局部紧空间，\(\mu\) 是 T 上的有界正测度，R 是 T 中的 \(\mu\) -可测等价关系。设 \(\Omega\) 是 T 上测度空间 \(\mathcal{M}(\mathrm{T})\) 中由满足 \(\lambda \geqslant 0\) 、总质量 \(\leqslant 1\) 且非零的测度组成的子集；\(\Omega\) 赋有淡拓扑（第 III 章 §1，第 9 号，命题 15 的推论 2）时是局部紧的。证明存在唯一的正测度 \(\rho\) （ \(\Omega\) 上的），使得： \(1^{\circ} \mu = \int_{\Omega} \lambda d\rho(\lambda)\) ； \(2^{\circ} \rho\) 集中于一个子集 \(B_{0}\) （ \(\Omega\) 中由总质量为 1、由模 R 的类承载的测度组成的子集）上，且 \(B_{0}\) 的两个不同元素由不同的类承载。（考虑分解 \(b \mapsto \lambda_{b}\) （ \(\mu\) 按关系 R 的一种分解），使所有 \(\lambda_{b}\) 都具有总质量 1，并考虑 B 在映射 \(b \mapsto \lambda_{b}\) （ B 到 \(\Omega\) 中的）下的像；利用定理 4。）
\item
  设 T 是 R 中的区间 \([0,2]\) ，\(\mu\) 是 T 上的 Lebesgue 测度。设 A 是含于 Cantor 集 \(K \subset [0,1]\) 的非 Borel 集（参见 GT，IV，§2，第 4 号，例以及 IX，§6，习题 6）。设 S 是 T 中的等价关系，其等价类是集合 \(\{x\}\) (对 \(x \notin A \cup (A + 1)\) )与集合 \(\{x, x + 1\}\) (对 \(x \in A\) )。证明 S 是 \(\mu\) -可测的，但不存在 S 的 Borel 截面。
\item
  设 \(\mathbf{K}\) 是可度量化紧空间，\(f\) 是 \(\mathbf{K}\) 到 Hausdorff 拓扑空间 \(\mathbf{E}\) 中的连续映射，\(k\) 是任意整数 \(> 1\) 。用 \(\mathrm{B}_k\) 表示 \(\mathbf{E}\) 中由这样的 \(y\) 组成的子集：\(\stackrel{-1}{f(y)}\) 至少含 \(k\) 个不同的点；证明 \(\mathrm{A}_k = \stackrel{-1}{f(\mathrm{B}_k)}\) 是 \(\mathbf{K}\) 中的 Borel 集（对每个整数 \(n > 0\) ，用 \(\mathrm{B}_{kn}\) 表示 \(\mathbf{E}\) 中由这样的 \(y\) 组成的子集：\(\stackrel{-1}{f(y)}\) 至少含 \(k\) 个彼此距离均为 \(\geqslant 1/n\) 的点；证明 \(\mathrm{B}_{kn}\) 是闭集）。
\end{enumerate}

¶ 10) 设 K 是可度量化紧空间，μ 是 K 上的正测度，f 是 K 到 Hausdorff 拓扑空间 E 中的 μ-可测映射。用 I（相应地 U）表示 E 中由这样的 y 组成的子集：\(\stackrel{-1}{f(y)}\) 是无限集（相应地不可数集）。

\begin{enumerate}
\def\labelenumi{\alph{enumi})}
\item
  证明 K 中存在 \(\mu\) -可测集 H，使得 \(f(\mathrm{H}) = \mathrm{I}\) ，且对每个 \(y \in \mathrm{E}\) ，\(\stackrel{-1}{f(y)} \cap \mathbb{C}\mathrm{H}\) 是有限集。（设 \((\mathrm{K}_n)\) 是 K 中紧子集的递增序列，使得 \(f\) 在每个 \(\mathrm{K}_n\) 上的限制连续，且诸 \(\mathrm{K}_n\) 之并 F 的余集 N 是 \(\mu\) -可忽略的；对每个 \(k > 1\) ，设 \(\mathrm{B}_k\) 是 E 中由这样的 \(y\) 组成的子集：\(\mathrm{F} \cap \stackrel{-1}{f(y)}\) 至少含 \(k\) 个不同的点，并设 \(\mathrm{A}_k = \mathrm{F} \cap \stackrel{-1}{f(\mathrm{B}_k)}\) ；取 H 为集合 \(\mathrm{A} = \bigcap_k \mathrm{A}_k\) 与集合 \(\mathrm{N} \cap \stackrel{-1}{f(\mathrm{I})}\) 之并，并利用习题 9。）
\item
  设 \(\mathfrak{F}\) 是由这样的 \(\mu\) -可测集 \(M \subset H\) 组成的集：\(M \cap f(y)\) 对每个 \(y \in E\) 都是有限集；设 \(\alpha\) 是测度 \(\mu(M)\) ( \(M \in \mathfrak{F}\) )的上确界；存在集合的递增序列 \((\mathrm{M}_n)\) （ \(\mathfrak{F}\) 中的），使得 \(\lim_{n\to \infty}\mu (\mathrm{M}_n) = \alpha\) 。设 \(\mathrm{P} = \bigcup_{n}\mathrm{M}_{n}\)
\end{enumerate}

并设 \(S\) 是 \(H \cap \mathbb{C}P\) 关于等价关系 \(f(x) = f(y)\) 的可测截面（定理 3）；证明 \(\mu(S) = 0\) 。

\begin{enumerate}
\def\labelenumi{\alph{enumi})}
\setcounter{enumi}{2}
\tightlist
\item
  证明存在 \(\mu\) -可忽略集 \(Z \subset K\) ，使得 \(f(Z) = U\) ，且对每个 \(\varepsilon > 0\) ，存在 \(\mu\) -可测集 \(L \subset K\) ，使得 \(\mu(L) \leqslant \varepsilon\) 且 \(f(L) = I\) （注意 \(f(H \cap \mathbb{C}M_n) = I\) 对每个 \(n\) 成立，且 \(U \subset f(S)\) ）。
\end{enumerate}

¶ 11) 设 K 是可度量化紧空间，μ 是 K 上的正测度，f 是 K 到局部紧空间 T 中的 μ-可测映射，ν 是 T 上的正测度；I 与 U 的含义同习题 10。

\begin{enumerate}
\def\labelenumi{\alph{enumi})}
\item
  假设对每个 \(\mu\) -可忽略集 \(N \subset K\) ，\(f(N)\) 都是 \(\nu\) -可忽略的。则 \(U\) 是 \(\nu\) -可忽略的，且对每个 \(\mu\) -可测子集 \(A \subset K\) ，\(f(A)\) 是 \(\nu\) -可测的。
\item
  为使 I 是 \(\nu\) -可忽略的且在 \(f\) 下每个 \(\mu\) -可忽略集的像是 \(\nu\) -可忽略的，必须且只需 \(f\) 满足以下条件：对每个 \(\varepsilon > 0\) ，存在 \(\delta > 0\) ，使得对每个子集 \(A \subset K\) ( \(\mu\) -可测且满足 \(\mu(A) \leqslant \delta\) 的)，\(f(A)\) 是 \(\nu\) -可测的且 \(\nu(f(A)) \leqslant \varepsilon\) 。（为证条件必要，用反证法，考虑递减序列 \((A_n)\) （由 \(\mu\) -可测集组成），其交是 \(\mu\) -可忽略的，但 \(\nu(f(A_n)) \geqslant \alpha > 0\) ；然后取诸 \(f(A_n)\) 与 \(\mathbb{C}\mathrm{I}\) 的交。）
\end{enumerate}

¶ 12) 设 E 是在 R 的区间 \([-1,+1]\) 上赋以 GT，IX，§2 的习题 13 a) 中定义的拓扑 T 而得到的非可度量化紧空间；设 \(\varphi\) 是映射 \(x \mapsto |x|\) （ E 到 I = {[}0,1{]} 上的，I 赋以 R 的拓扑所诱导的拓扑）。

\begin{enumerate}
\def\labelenumi{\alph{enumi})}
\item
  证明 \(\varphi\) 连续，并且，若用 \(\mathfrak{F}\) 表示 \(E\) 的形如 \(\overline{\varphi}(A)\) 的子集组成的集（其中 \(A\) 是 \(I\) 的 Borel 子集），则 \(E\) 的 Borel 子集是这样的子集 \(B\) ：存在 \(M \in \mathfrak{F}\) ，使得 \(B \cap \mathbb{C}M\) 与 \(M \cap \mathbb{C}B\) 均可数（注意 \(E\) 的每个开集都属于 \(\mathfrak{F}\) ）。
\item
  证明数值函数 \(f\) （定义在 \(\mathbf{E}\) 上的）为连续（对 \(\mathcal{T}\) ），必须且只需它是正则函数，且 \(f(x-) = f((-x)+)\) 对每个 \(x \in \mathbf{E}\) 成立。于是定义正测度 \(\mu\) （ \(\mathbf{E}\) 上的）如下：令 \(\int f d\mu = \int_0^1 f(x) dx\) 。证明 \(\mu\) 在 \(\varphi\) 下的像是 \(\mathbf{I}\) 上的 Lebesgue 测度，且 \(\mu\) -可忽略集正是 \([-1, +1]\) 上对 Lebesgue 测度可忽略的子集。
\item
  由 \(a)\) 与 \(b)\) 推出：不存在 \(\mu\) -可测截面，对应于 \(\mu\) -可测等价关系 \(\varphi(x) = \varphi(y)\) （在 \(E\) 中）。
\end{enumerate}

